% Multi-Source Knowledge Graph RAG for Medical Diagnosis
% v17 - 12-page condensed version
% Target: Information Fusion (Elsevier)

\documentclass[preprint,12pt]{elsarticle}

\usepackage{amsmath,amssymb,amsfonts}
\usepackage{graphicx}
\usepackage{booktabs}
\usepackage{multirow}
\usepackage{hyperref}
\usepackage{algorithm}
\usepackage{algorithmic}
\usepackage{xcolor}
\usepackage{url}

\journal{Information Fusion}

\begin{document}

\begin{frontmatter}

\title{Multi-Source Knowledge Graph Retrieval-Augmented Generation\\as a Second Brain for LLM-Based Medical Diagnosis}

\author[addr1]{First Author}
\author[addr1]{Second Author}
\author[addr1]{Third Author}

\address[addr1]{Department of Computer Science, University, Country}

\begin{abstract}
Large Language Models (LLMs) suffer from hallucination and knowledge gaps in medical domains. We propose a multi-source Retrieval-Augmented Generation (RAG) framework functioning as a ``Second Brain'' for LLM-based differential diagnosis, integrating three complementary retrieval layers: dense vector retrieval, sparse BM25 retrieval enhanced with Positive Pointwise Mutual Information (PPMI) co-occurrence statistics, and a multi-source knowledge graph combining ontological structure with data-driven co-occurrence patterns. On DDXPlus (49 diseases, $n{=}200$), our pipeline achieves 95.5\% Top-1 and 98.5\% Top-5 accuracy with zero hallucinations, versus 32.5\% Top-1 and 59.8\% hallucination for the unaugmented LLM. On Symptom2Disease (24 diseases, $n{=}320$), it achieves 90.3\% Top-1 and 98.8\% Top-5 with zero hallucinations. A GPT-4o ablation shows both models converge to identical performance under RAG, demonstrating that retrieval compensates for base model variation.
\end{abstract}

\begin{keyword}
Retrieval-Augmented Generation \sep Knowledge Graph \sep Medical Diagnosis \sep Large Language Models \sep Hallucination Mitigation \sep Second Brain
\end{keyword}

\end{frontmatter}

%% ============================================================
\section{Introduction}
\label{sec:introduction}

Large Language Models (LLMs) have demonstrated impressive capabilities across diverse domains~\cite{brown2020language,openai2023gpt4}, yet deploying them in safety-critical medical applications reveals fundamental limitations: hallucination of plausible but incorrect facts, lack of grounding in verified clinical knowledge, and inability to reliably perform structured differential diagnosis~\cite{ji2023survey,singhal2023large}.

Retrieval-Augmented Generation (RAG) addresses these limitations by grounding LLM outputs in retrieved evidence~\cite{lewis2020retrieval,gao2023retrieval}. However, standard dense retrieval faces challenges in medical domains: semantic similarity in embedding space does not always capture clinical relevance, and single-source retrieval may miss complementary signals. We propose a \textbf{multi-source RAG framework} that functions as a ``Second Brain'' for LLMs performing medical diagnosis, inspired by Forte's methodology~\cite{forte2022building} of capturing and retrieving knowledge through complementary systems. Our architecture integrates: (1)~dense vector retrieval capturing semantic relationships, (2)~sparse BM25 retrieval enhanced with PPMI co-occurrence statistics capturing statistical associations, and (3)~a multi-source knowledge graph combining ontological structure with data-driven co-occurrence patterns.

Different retrieval modalities capture complementary aspects of medical knowledge. Dense embeddings excel at semantic similarity but may miss statistically significant associations. Sparse retrieval with co-occurrence statistics captures these patterns but lacks structural context. Knowledge graphs provide structural relationships but require careful construction. By fusing all three, we achieve near-perfect diagnostic accuracy with zero hallucination.

Our contributions are: (1)~a principled multi-source RAG architecture progressively integrating dense, sparse, and graph-based retrieval (\S\ref{sec:method}); (2)~empirical demonstration on DDXPlus~\cite{fansi2022ddxplus} and Symptom2Disease~\cite{gretelai_s2d} achieving 98.5\% and 98.8\% Top-5 accuracy (\S\ref{sec:results}); (3)~evidence that multi-source RAG completely eliminates hallucination (\S\ref{sec:hallucination}); and (4)~a model ablation showing RAG compensates for base model differences (\S\ref{sec:model_ablation}).

%% ============================================================
\section{Related Work}
\label{sec:related}

\textbf{Retrieval-Augmented Generation.} RAG was introduced by Lewis et al.~\cite{lewis2020retrieval} to augment language models with retrieved documents, combining dense passage retrieval~\cite{karpukhin2020dense} with sequence-to-sequence generation. It has since been extended to open-domain QA~\cite{izacard2021leveraging}, dialogue~\cite{shuster2021retrieval}, and fact verification~\cite{thorne2018fever}. Recent surveys~\cite{gao2023retrieval,zhao2024retrieval} categorize approaches by retrieval granularity, fusion strategy, and retrieval-generation coupling. Our work contributes to the \emph{multi-source retrieval} dimension, fusing heterogeneous dense, sparse, and graph-based signals.

\textbf{Knowledge Graphs in Biomedical AI.} Knowledge graphs structure biomedical knowledge through resources such as UMLS~\cite{bodenreider2004unified} and domain-specific ontologies~\cite{santos2022knowledge,chandak2023building}. Recent work integrates KGs with LLMs for medical reasoning~\cite{pan2024unifying}, including graph-enhanced retrieval~\cite{he2024gretriever} and knowledge-grounded generation~\cite{yasunaga2022deep}. Unlike prior approaches using KGs as static lookup tables, we construct a \emph{multi-source} KG combining ontological edges with data-driven co-occurrence and PPMI-weighted edges as one component of a broader retrieval pipeline.

\textbf{LLM-Based Medical Diagnosis and Hallucination.} Med-PaLM~\cite{singhal2023large} showed LLMs can approach physician-level medical QA performance, but differential diagnosis remains challenging~\cite{mcduff2023towards}. Hallucination---generating fluent but factually incorrect content---is particularly dangerous in medical contexts~\cite{ji2023survey,umapathi2023med}. Mitigation strategies include retrieval augmentation~\cite{shuster2021retrieval}, constrained decoding~\cite{lu2022neurologic}, and self-consistency checking~\cite{manakul2023selfcheckgpt}. We demonstrate that multi-source RAG can \emph{completely eliminate} hallucination in structured diagnostic tasks.

%% ============================================================
\section{Method}
\label{sec:method}

\subsection{Problem Formulation}

Given patient symptoms $S = \{s_1, s_2, \ldots, s_m\}$ and optional demographics, the task is to produce a ranked diagnosis list $\hat{D} = [\hat{d}_1, \ldots, \hat{d}_k]$ where ground-truth $d^*$ should rank as high as possible ($k{=}5$). We measure Top-$k$ Accuracy ($\mathbb{1}[d^* \in \hat{D}_{1:k}]$), NDCG@5, Macro F1, and Hallucination Rate ($\frac{1}{N}\sum_{i=1}^{N} \mathbb{1}[\exists \hat{d} \in \hat{D}_i : \hat{d} \notin \mathcal{D}]$).

\subsection{Architecture Overview}

Our framework processes each case through progressive retrieval modules fused before LLM prompting. We evaluate four configurations: (1)~\textbf{LLM-Only}: direct prompting; (2)~\textbf{Dense RAG}: adding dense vector retrieval; (3)~\textbf{BM25+PPMI}: adding sparse retrieval with co-occurrence statistics; (4)~\textbf{Multi-Source KG}: adding knowledge graph traversal. Figure~\ref{fig:architecture} shows the architecture.

\begin{figure}[t]
\centering
\fbox{\parbox{0.9\linewidth}{\centering
\textbf{Multi-Source RAG Architecture}\\[4pt]
Patient Symptoms $\rightarrow$ \{Dense Retrieval $\|$ BM25+PPMI $\|$ KG Traversal\}\\
$\downarrow$\\
Reciprocal Rank Fusion $\rightarrow$ Augmented Prompt $\rightarrow$ LLM $\rightarrow$ Ranked Diagnoses
}}
\caption{Multi-source RAG pipeline overview.}
\label{fig:architecture}
\end{figure}

\subsection{Dense Vector Retrieval}

Each disease $d_j$ is represented by a profile document (name, symptoms with prevalence, description) embedded using OpenAI's \texttt{text-embedding-3-small} (1536 dimensions). At inference, the symptom set is embedded as a query and the top-10 most similar profiles retrieved via cosine similarity:
\begin{equation}
    \text{sim}(q, d_j) = \frac{\mathbf{e}_q \cdot \mathbf{e}_{d_j}}{\|\mathbf{e}_q\| \cdot \|\mathbf{e}_{d_j}\|}
\end{equation}

\subsection{BM25 with PPMI Enhancement}

Disease profiles are indexed with BM25~\cite{robertson2009probabilistic}, capturing lexical overlap missed by dense embeddings. We enhance retrieval with PPMI co-occurrence scores from training data:
\begin{equation}
    \text{PPMI}(s_i, d_j) = \max\left(0, \log_2 \frac{P(s_i, d_j)}{P(s_i) \cdot P(d_j)}\right)
\end{equation}
The combined score is $\text{score}_{\text{BM25+PPMI}}(d_j | S) = \text{BM25}(d_j | S) + \lambda \sum_{s_i \in S} \text{PPMI}(s_i, d_j)$ with $\lambda = 0.5$.

\subsection{Multi-Source Knowledge Graph}

We construct a heterogeneous graph $G = (V, E)$ with disease and symptom nodes connected by three edge types: \emph{ontological} ($E_O$, binary associations from disease definitions), \emph{co-occurrence} ($E_C$, weighted by normalized frequency), and \emph{PPMI} ($E_P$, weighted by PPMI scores). Given query symptoms $S$, we traverse from symptom nodes to neighboring disease nodes with aggregated score:
\begin{equation}
    \text{score}_{\text{KG}}(d_j | S) = \sum_{s_i \in S} \left(\alpha \cdot w_O(s_i, d_j) + \beta \cdot w_C(s_i, d_j) + \gamma \cdot w_P(s_i, d_j)\right)
\end{equation}
where $\alpha = \beta = \gamma = 1/3$.

\subsection{Retrieval Fusion and LLM Prompting}

Retrieved candidates are merged via Reciprocal Rank Fusion~\cite{cormack2009reciprocal}: $\text{RRF}(d_j) = \sum_{r \in R} \frac{1}{60 + \text{rank}_r(d_j)}$. The top-10 disease profiles are formatted into structured context included in the LLM prompt, which instructs the model to consider only provided candidates and output ranked diagnoses as JSON, directly mitigating hallucination through constrained generation.

%% ============================================================
\section{Experimental Setup}
\label{sec:experiments}

\textbf{DDXPlus~\cite{fansi2022ddxplus}} contains ${\sim}$1.3M patient cases across 49 pathologies and 223 symptoms, constructed from a validated clinical knowledge base. We evaluate on $n{=}200$ stratified samples from the 134K test set. The KG and PPMI statistics are built from 10K training cases.

\textbf{Symptom2Disease~\cite{gretelai_s2d}} contains ${\sim}$1,065 free-text symptom descriptions mapped to 24 diseases. We split 70/30 into training ($n{=}745$) and test ($n{=}320$), constructing a \emph{domain-specific} KG entirely from S2D training data to test cross-domain generalization.

\textbf{Models.} GPT-4o-mini (\texttt{gpt-4o-mini-2024-07-18}) serves as primary LLM with temperature 0.0 and seed 42. GPT-4o (\texttt{gpt-4o-2024-08-06}) is used for model ablation. Dense embeddings use \texttt{text-embedding-3-small}. Hallucination is measured by checking whether each predicted disease appears in the target dataset's vocabulary.

%% ============================================================
\section{Results}
\label{sec:results}

\subsection{DDXPlus Results and Model Ablation}
\label{sec:model_ablation}

Table~\ref{tab:ddxplus_combined} presents results on DDXPlus for both GPT-4o-mini and GPT-4o. The unaugmented GPT-4o-mini achieves only 32.5\% Top-1 with 59.8\% hallucination. Dense RAG improves to 85.5\% Top-1 and eliminates hallucination. BM25+PPMI further improves Top-1 to 95.5\%---a 10pp gain---demonstrating that sparse retrieval captures complementary signals. The full Multi-Source KG matches 95.5\% Top-1 but improves Top-3 to 98.5\% and NDCG@5 to 0.999, primarily enhancing \emph{ranking quality}.

\begin{table}[t]
\centering
{\small
\caption{Results on DDXPlus ($n{=}200$, 49 diseases). GPT-4o-mini (top) and GPT-4o ablation (bottom). Best per model in \textbf{bold}.}
\label{tab:ddxplus_combined}
\begin{tabular}{llcccccc}
\toprule
\textbf{Model} & \textbf{Method} & \textbf{Top-1} & \textbf{Top-3} & \textbf{Top-5} & \textbf{NDCG@5} & \textbf{F1} & \textbf{Halluc.} \\
\midrule
\multirow{4}{*}{\rotatebox{90}{\scriptsize 4o-mini}}
& LLM-Only       & 0.325 & 0.420 & 0.485 & 0.415 & 0.163 & 0.598 \\
& Dense RAG       & 0.855 & 0.935 & 0.955 & 0.934 & 0.321 & \textbf{0.000} \\
& BM25+PPMI       & 0.955 & 0.965 & 0.965 & 0.993 & 0.325 & 0.008 \\
& Multi-Source KG & \textbf{0.955} & \textbf{0.985} & \textbf{0.985} & \textbf{0.999} & \textbf{0.331} & \textbf{0.000} \\
\midrule
\multirow{4}{*}{\rotatebox{90}{\scriptsize GPT-4o}}
& LLM-Only       & 0.235 & 0.365 & 0.430 & 0.356 & 0.151 & 0.672 \\
& Dense RAG       & 0.820 & 0.925 & 0.955 & 0.918 & 0.321 & 0.002 \\
& BM25+PPMI       & 0.960 & 0.970 & 0.970 & 0.992 & 0.326 & 0.013 \\
& Multi-Source KG & \textbf{0.955} & \textbf{0.985} & \textbf{0.985} & \textbf{0.999} & \textbf{0.331} & 0.003 \\
\bottomrule
\end{tabular}
}
\end{table}

A striking finding emerges from the GPT-4o ablation: the larger model performs \emph{worse} without RAG (23.5\% vs.\ 32.5\% Top-1, 67.2\% vs.\ 59.8\% hallucination), suggesting its broader knowledge introduces more spurious associations when unconstrained. However, with full RAG augmentation, both models converge to identical performance (95.5\% Top-1, 98.5\% Top-5, NDCG@5 of 0.999). This convergence validates that \textbf{our retrieval framework compensates for variation in base model capability}, implying that investment in retrieval infrastructure yields greater returns than larger base models for structured diagnostic tasks.

The incremental contribution of each layer on DDXPlus follows a clear pattern: Dense RAG provides +53.0pp Top-1 (the foundation), BM25+PPMI adds +10.0pp by capturing co-occurrence patterns missed by embeddings, and the KG adds +2.0pp Top-5 while optimizing NDCG@5 from 0.993 to 0.999 and eliminating residual hallucination.

\subsection{Cross-Domain Results: Symptom2Disease}
\label{sec:cross_domain}

Table~\ref{tab:s2d} shows results on S2D with a domain-specific KG built from S2D training data.

\begin{table}[t]
\centering
{\small
\caption{Results on Symptom2Disease ($n{=}320$, 24 diseases) using GPT-4o-mini with domain-specific KG.}
\label{tab:s2d}
\begin{tabular}{lcccccc}
\toprule
\textbf{Method} & \textbf{Top-1} & \textbf{Top-3} & \textbf{Top-5} & \textbf{NDCG@5} & \textbf{F1} & \textbf{Halluc.} \\
\midrule
LLM-Only       & 0.325 & 0.478 & 0.569 & 0.494 & 0.196 & 0.759 \\
Dense RAG       & 0.800 & 0.916 & 0.953 & 0.883 & 0.318 & \textbf{0.000} \\
BM25+PPMI       & 0.766 & 0.794 & 0.794 & 0.784 & 0.265 & 0.004 \\
Multi-Source KG & \textbf{0.903} & \textbf{0.978} & \textbf{0.988} & \textbf{0.952} & \textbf{0.329} & \textbf{0.000} \\
\bottomrule
\end{tabular}
}
\end{table}

The pipeline generalizes effectively: Multi-Source KG achieves 90.3\% Top-1 and 98.8\% Top-5 with zero hallucinations, versus 32.5\% Top-1 and 75.9\% hallucination for LLM-Only. Notably, BM25+PPMI \emph{underperforms} Dense RAG on S2D (76.6\% vs.\ 80.0\% Top-1)---a reversal from DDXPlus---reflecting that BM25's lexical matching is less effective for S2D's varied free-text descriptions compared to DDXPlus's structured symptom codes. The Multi-Source KG provides its largest relative improvement here (+10.3pp over Dense RAG), validating the multi-source design: when individual modalities show mixed performance, the KG's integration of all signals yields the greatest benefit.

\subsection{Hallucination Analysis}
\label{sec:hallucination}

Table~\ref{tab:halluc} summarizes hallucination rates across all conditions.

\begin{table}[t]
\centering
{\small
\caption{Hallucination rates across methods, datasets, and models.}
\label{tab:halluc}
\begin{tabular}{lccc}
\toprule
\textbf{Method} & \textbf{DDXPlus (4o-mini)} & \textbf{DDXPlus (4o)} & \textbf{S2D (4o-mini)} \\
\midrule
LLM-Only        & 0.598 & 0.672 & 0.759 \\
Dense RAG        & 0.000 & 0.002 & 0.000 \\
BM25+PPMI        & 0.008 & 0.013 & 0.004 \\
Multi-Source KG  & 0.000 & 0.003 & 0.000 \\
\bottomrule
\end{tabular}
}
\end{table}

Without RAG, hallucination rates reach 60--76\%, with GPT-4o \emph{higher} than GPT-4o-mini (67.2\% vs.\ 59.8\%), suggesting larger models generate more diverse but invalid disease names. S2D shows the highest baseline hallucination (75.9\%), reflecting the model's vocabulary exceeding the 24-disease scope. The full pipeline eliminates hallucination entirely with GPT-4o-mini and reduces it to negligible 0.3\% with GPT-4o.

%% ============================================================
\section{Discussion}
\label{sec:discussion}

\textbf{Progressive value of multi-source retrieval.} Each retrieval layer captures a distinct aspect of disease-symptom relationships: dense embeddings capture semantic similarity, BM25+PPMI captures statistical co-occurrence patterns, and the KG captures structural ontological relationships weighted by multiple evidence sources. This complementarity mirrors clinical reasoning---combining textbook knowledge, clinical experience, and pattern matching. The S2D results further illuminate this interplay: when sparse retrieval underperforms dense retrieval (due to free-text vs.\ structured inputs), the KG compensates by integrating both signals, achieving best performance regardless of input format.

\textbf{RAG as a model equalizer.} The convergence of GPT-4o and GPT-4o-mini under full RAG (both at 95.5\% Top-1, 98.5\% Top-5) has significant practical implications. The base model's parametric knowledge matters less than retrieval quality; smaller, cheaper models match larger ones with appropriate retrieval infrastructure; and the retrieval framework provides a performance floor largely independent of model size. Organizations can deploy cost-effective models without sacrificing accuracy.

\textbf{Domain-specific knowledge graphs.} The strong S2D results demonstrate generalization to new medical domains via automatic KG construction from training data, requiring no expert curation. The performance gap between DDXPlus (95.5\%) and S2D (90.3\%) reflects genuine difficulty differences: S2D's free-text descriptions have higher linguistic variability and fewer training examples per disease (${\sim}$31 vs.\ ${\sim}$204). Hybrid approaches combining automatic construction with expert curation could further improve performance.

\textbf{Clinical significance proxies.} Our use of TF-IDF and PPMI effectively captures population-level patterns but may not reflect pathognomonic symptoms (rare but clinically definitive), severity weighting, or temporal patterns. Integrating expert-curated specificity weights from clinical knowledge bases could address these limitations.

\textbf{Practical deployment.} Near-perfect performance (98.5--98.8\% Top-5, zero hallucination) suggests viability as clinical decision support within well-defined domains, with considerations for knowledge base maintenance, scope transparency, domain adaptation via automatic KG construction, and regulatory compliance.

%% ============================================================
\section{Limitations}
\label{sec:limitations}

Our evaluation covers 49 (DDXPlus) and 24 (S2D) diseases---subsets of real clinical pathology---with DDXPlus cases synthetically generated and potentially missing comorbidities and atypical presentations. Sample sizes ($n{=}200$ and $n{=}320$) suffice for large effect sizes but warrant larger-scale validation. We evaluate only GPT-4o-mini and GPT-4o; interactions with open-source medical models (e.g., Meditron~\cite{chen2023meditron}) may differ. Our static KGs do not incorporate external resources (UMLS, SNOMED-CT) or adapt over time, and we use equal fusion weights ($\alpha{=}\beta{=}\gamma{=}1/3$) rather than learned weights---though near-perfect results leave limited room for tuning.

%% ============================================================
\section{Conclusion}
\label{sec:conclusion}

We presented a multi-source RAG framework serving as a ``Second Brain'' for LLM-based medical diagnosis. By progressively integrating dense vector retrieval, BM25 with PPMI co-occurrence statistics, and a multi-source knowledge graph, our system achieves 95.5\% Top-1 and 98.5\% Top-5 accuracy on DDXPlus and 90.3\% Top-1 and 98.8\% Top-5 on Symptom2Disease---with zero hallucinations---compared to unaugmented baselines of 32.5\% Top-1 and 60--76\% hallucination. Each layer contributes complementary information: dense retrieval provides the largest gain, BM25+PPMI captures statistical patterns missed by embeddings, and the KG optimizes ranking while eliminating residual hallucinations. A GPT-4o ablation confirms that retrieval compensates for base model variation, with both models converging under full augmentation. The pipeline generalizes to new domains via automatic KG construction, offering a principled path toward reliable clinical decision support.

%% ============================================================
\section*{CRediT Authorship Contribution Statement}

\textbf{First Author:} Conceptualization, Methodology, Software, Validation, Writing -- Original Draft.
\textbf{Second Author:} Methodology, Formal Analysis, Writing -- Review \& Editing.
\textbf{Third Author:} Supervision, Writing -- Review \& Editing.

\section*{Declaration of Competing Interest}

The authors declare no known competing financial interests or personal relationships that could have influenced this work.

\section*{Data Availability}

DDXPlus: \url{https://github.com/mila-iqia/ddxplus}. Symptom2Disease: \url{https://huggingface.co/datasets/gretelai/symptom_to_diagnosis}. Code will be released upon publication.

%% ============================================================
\bibliographystyle{elsarticle-num}

\begin{thebibliography}{27}

\bibitem{brown2020language}
T.~Brown, B.~Mann, N.~Ryder, et~al., Language models are few-shot learners, in: NeurIPS, vol.~33, 2020, pp.~1877--1901.

\bibitem{openai2023gpt4}
OpenAI, GPT-4 technical report, arXiv:2303.08774, 2023.

\bibitem{ji2023survey}
Z.~Ji, N.~Lee, R.~Frieske, et~al., Survey of hallucination in natural language generation, ACM Comput. Surv. 55~(12) (2023) 1--38.

\bibitem{singhal2023large}
K.~Singhal, S.~Azizi, T.~Tu, et~al., Large language models encode clinical knowledge, Nature 620~(7972) (2023) 172--180.

\bibitem{lewis2020retrieval}
P.~Lewis, E.~Perez, A.~Piktus, et~al., Retrieval-augmented generation for knowledge-intensive NLP tasks, in: NeurIPS, vol.~33, 2020, pp.~9459--9474.

\bibitem{gao2023retrieval}
Y.~Gao, Y.~Xiong, A.~Gupta, et~al., Retrieval-augmented generation for large language models: A survey, arXiv:2312.10997, 2023.

\bibitem{forte2022building}
T.~Forte, Building a Second Brain, Atria Books, 2022.

\bibitem{robertson2009probabilistic}
S.~Robertson, H.~Zaragoza, The probabilistic relevance framework: BM25 and beyond, Found. Trends Inf. Retr. 3~(4) (2009) 333--389.

\bibitem{fansi2022ddxplus}
A.~Fansi~Tchango, R.~Camara, I.~Kreitchmann, et~al., DDXPlus: A new dataset for automatic medical diagnosis, in: NeurIPS, vol.~35, 2022.

\bibitem{gretelai_s2d}
Gretel.ai, Symptom to diagnosis dataset, HuggingFace, 2023. URL: \url{https://huggingface.co/datasets/gretelai/symptom_to_diagnosis}.

\bibitem{santos2022knowledge}
A.~Santos, A.~Coelho, et~al., A knowledge graph to interpret clinical proteomics data, Nat. Biotechnol. 40 (2022) 692--702.

\bibitem{chandak2023building}
P.~Chandak, K.~Huang, M.~Zitnik, Building a knowledge graph to enable precision medicine, Sci. Data 10 (2023) 67.

\bibitem{bodenreider2004unified}
O.~Bodenreider, The UMLS: Integrating biomedical terminology, Nucleic Acids Res. 32 (2004) D267--D270.

\bibitem{pan2024unifying}
S.~Pan, L.~Luo, Y.~Wang, et~al., Unifying large language models and knowledge graphs: A roadmap, IEEE TKDE (2024).

\bibitem{mcduff2023towards}
D.~McDuff, M.~Schaekermann, T.~Tu, et~al., Towards accurate differential diagnosis with large language models, arXiv:2312.00164, 2023.

\bibitem{karpukhin2020dense}
V.~Karpukhin, B.~O\u{g}uz, S.~Min, et~al., Dense passage retrieval for open-domain question answering, in: EMNLP, 2020, pp.~6769--6781.

\bibitem{izacard2021leveraging}
G.~Izacard, E.~Grave, Leveraging passage retrieval with generative models for open domain question answering, in: EACL, 2021, pp.~874--880.

\bibitem{shuster2021retrieval}
K.~Shuster, S.~Poff, M.~Chen, et~al., Retrieval augmentation reduces hallucination in conversation, in: EMNLP Findings, 2021.

\bibitem{thorne2018fever}
J.~Thorne, A.~Vlachos, C.~Christodoulopoulos, A.~Mittal, FEVER: A large-scale dataset for fact extraction and verification, in: NAACL-HLT, 2018, pp.~809--819.

\bibitem{zhao2024retrieval}
P.~Zhao, H.~Zhang, Q.~Yu, et~al., Retrieval-augmented generation for AI-generated content: A survey, arXiv:2402.19473, 2024.

\bibitem{umapathi2023med}
L.~K.~Umapathi, A.~Pal, M.~Sankarasubbu, Med-HALT: Medical domain hallucination test for large language models, in: CoNLL, 2023.

\bibitem{lu2022neurologic}
X.~Lu, S.~Welleck, P.~West, et~al., NeuroLogic A*esque decoding, in: NAACL, 2022.

\bibitem{manakul2023selfcheckgpt}
P.~Manakul, A.~Liusie, M.~J.~F.~Gales, SelfCheckGPT: Zero-resource black-box hallucination detection, in: EMNLP, 2023.

\bibitem{cormack2009reciprocal}
G.~V.~Cormack, C.~L.~A.~Clarke, S.~B\"uttcher, Reciprocal rank fusion outperforms condorcet and individual rank learning methods, in: SIGIR, 2009, pp.~758--759.

\bibitem{he2024gretriever}
X.~He, Y.~Tian, Y.~Sun, et~al., G-Retriever: Retrieval-augmented generation for textual graph understanding, arXiv:2402.07630, 2024.

\bibitem{yasunaga2022deep}
M.~Yasunaga, A.~Bosselut, H.~Ren, et~al., Deep bidirectional language-knowledge graph pretraining, in: NeurIPS, 2022.

\bibitem{chen2023meditron}
Z.~Chen, A.~Cano, et~al., Meditron-70B: Scaling medical pretraining for large language models, arXiv:2311.16079, 2023.

\end{thebibliography}

\end{document}
